% Checkpoint36 proposal. No checkpoint35 formal source is changed.
\statement{A.124}{Proved}{The actual elliptic cubic has a unique continuous pure transporter and an exact involution obstruction}
Let \(\mathcal G=B_{1,3}\) be the discrete reduced elliptic braid
group in A.121, and put
\(\mathcal P=P_{1,3}=\ker(\mathcal G\to S_3)\).
All completions in this statement are full profinite completions.
Use the actual marked generators \(\sigma,\tau,X,Y\), and put
\[
 L=\tau^{-1}\sigma^{-1},\qquad \Pi_0=LY,
 \qquad \bar r=\pi(\Pi_0)\in S_3.
 \tag{A.124.1}
\]
There are exactly four conjugacy classes of nontrivial finite subgroups
of \(\widehat{\mathcal G}\), all cyclic of order three, and eight
conjugacy classes of nonidentity torsion elements. The latter are
distinguished exactly by their eight nonzero images in the indicated
odd factor of
\[
 \mathcal G^{\rm ab}
   =C_2\langle\sigma\rangle\times
     C_3\langle X\rangle\times C_3\langle Y\rangle.
 \tag{A.124.2}
\]
The centralizer and normalizer of every nontrivial finite subgroup
of \(\widehat{\mathcal G}\) equal that subgroup itself.

For every profinite commutator word \(w\) satisfying A.121's original
pentagon, let \(\Pi_w\) be its specified elliptic cubic. There is a
unique \(p_w\in\widehat{\mathcal P}\) such that
\[
                  \Pi_w=p_w\Pi_0p_w^{-1}.
 \tag{A.124.3}
\]
The map \(w\mapsto p_w\) is continuous on the closed space of these
words, and \(p_1=1\). In particular this applies to the actual canonical
word \(u\) of the original negative full-GT involution \(\alpha\).

Let \(\Theta\) be its actual elliptic involution, as in A.122.
There is a unique \(h_u\in\widehat{\mathcal P}\) with
\[
 \Theta(\Pi_u)=h_u\Pi_uh_u^{-1},\qquad
 \Theta(h_u)h_u=1.
 \tag{A.124.4}
\]
Write \(z_u=h_u^{-1}\) and let \(\vartheta\) denote the order-two
generator in
\(\widehat{\mathcal G}\rtimes_\Theta\langle\vartheta\rangle\).
Then \(z_u\vartheta\) is an involution commuting with \(\Pi_u\), and
\[
              \langle\Pi_u,z_u\vartheta\rangle\cong C_6.
 \tag{A.124.5}
\]
Its section class \([z_u]\in H^1(C_2,\widehat{\mathcal P};\Theta)\)
is trivial exactly when a pure conjugate of \(\Pi_u\) is fixed by
\(\Theta\). This class is the same for every \(\Pi_w\) above, with
\(\Theta\) held fixed. Neither its triviality nor the original
third-cusp equation or CM parity is asserted.

\begin{proof}
\emph{A.124.1. The precise group and its completion.}
Write \(T=\mathbb R^2/\mathbb Z^2\) and
\[
 \mathcal X=\{(a,b)\in T^2:a\ne0,\ b\ne0,\ a\ne b\}.
 \tag{A.124.6}
\]
Sending \((t_1,t_2,t_3)\) to
\((t_1,t_2-t_1,t_3-t_1)\) identifies the ordered configuration
space with \(T\times\mathcal X\).
The symmetric group acts by permuting the three points and then
renormalizing their first coordinate to zero.
The group \(\mathcal G\) is the group of lifts of this \(S_3\)-action
to the universal cover \(\widetilde{\mathcal X}\).
Equivalently, it is the \emph{orbifold} fundamental group of
\([\mathcal X/S_3]\). It is not the fundamental group of the coarse
topological quotient. Indeed the ordinary torus braid group, divided
by the two diagonal translation loops, gives this lift group and the
exact sequence
\[
       1\longrightarrow\pi_1(\mathcal X)
        \longrightarrow\mathcal G\longrightarrow S_3
        \longrightarrow1.
 \tag{A.124.7}
\]
The diagonal translation loops are the two total strand products;
killing them gives precisely the reduced presentation of
\cite[Section~1.2, equation~(7)]{A_Enriquez2014} used in A.121.
Thus (A.124.7) is the actual marked group, with
\(\mathcal P=\pi_1(\mathcal X)\).

Forgetting \(b\) gives a locally trivial fibration over \(T-\{0\}\),
with fiber \(T-\{0,a\}\). Local trivializations move \(a\) in a small
disk disjoint from zero by an ambient isotopy. The base and fiber
are aspherical, with free fundamental groups of ranks two and three.
The homotopy exact sequence and the CW structure therefore show that
\(\widetilde{\mathcal X}\) is contractible. The fundamental-group
sequence splits because its quotient is free, so
\(\mathcal P\cong F_3\rtimes F_2\).

We verify the full topology and cohomology needed below.
Given a finite quotient of \(F_3\), choose a characteristic finite-index
subgroup inside its kernel. The action of \(F_2\) on the resulting
finite group has finite image. Its finite semidirect product extends
the prescribed separation on \(F_3\). Together with finite quotients
of \(F_2\), this proves residual finiteness and the full induced
topology on the actual normal \(F_3\). The split sequence remains
exact on full completions. Comparing its discrete and profinite
Lyndon--Hochschild--Serre spectral sequences proves goodness of
\(\mathcal P\): the two free groups are good, and their finite-module
cohomology is finite in each degree and zero above degree one.
The same comparison for the finite quotient in (A.124.7) proves
goodness of \(\mathcal G\).

For completeness these arguments cover all discrete modules in the
hypothesis of the theorem we use. A finite graph model for \(F_3\),
with mapping cylinders for representatives of the two automorphisms
of its free quotient, gives a finite aspherical graph of spaces for
\(\mathcal P\). Its cellular resolution and the finite-group bar
resolution give resolutions of finite rank in each degree for the
cohomology comparisons. They commute with filtered unions. Every
discrete continuous \(\mathbb F_p[[\widehat{\mathcal G}]]\)-module
is the union of finite invariant submodules, so the finite-module
comparison extends to all such modules. Residual finiteness of
\(\mathcal G\) follows by taking characteristic finite-index
refinements in the finitely generated residually finite group
\(\mathcal P\); they are normal in \(\mathcal G\).
Finally \(\mathcal P\) is torsion-free of cohomological dimension
at most two, and the same spectral sequences give finite cohomology
for finite modules in every degree. Hence \(\mathcal G\) has finite
virtual \(p\)-cohomological type for every prime \(p\).

Theorem~A(i)--(iii) of \cite{BoggiZalesskii2025} therefore applies
\emph{directly} to \(\mathcal G\) and its full completion. It
preserves the conjugacy classes of finite prime-power subgroups and
of their elements, and identifies their completed centralizers and
normalizers with the closures of the discrete ones. We use its
prime-order case; no centralizer theorem for a nonsolvable subgroup
or completion identification inside a mapping class group is involved.

\emph{A.124.2. Discrete torsion and its unique fixed point.}
The pure group \(F_3\rtimes F_2\) is torsion-free. Thus a subgroup
of order two or three in \(\mathcal G\) injects into \(S_3\).
A three-cycle on \(\mathcal X\) is
\((a,b)\mapsto(b-a,-a)\). Its fixed points are precisely
\[
           (a,2a),\qquad a\in T[3]-\{0\}.
 \tag{A.124.8}
\]
There are eight, and they are isolated. A transposition has no
fixed point: in a configuration invariant up to translation, its
fixed marked point forces that translation to be zero, after which
its exchanged pair would coincide.

Here is a full fixed-point argument for a finite cyclic subgroup
\(C\) of prime order \(p=2\) or \(3\) acting on the actual cover.
Choose an \(S_3\)-invariant locally finite triangulation of
\(\mathcal X\), including (A.124.8) as vertices, and take a barycentric
subdivision to remove inversions. Such a triangulation can be made
explicitly: triangulate the compact torus \(T^4\) with the three
rational subtori \(a=0,b=0,a=b\) as subpolyhedra, take a common
refinement with its finitely many translates, and triangulate the
complement by layers of an invariant piecewise linear distance to
those subpolyhedra. Lift this triangulation to
\(\widetilde{\mathcal X}\). Its dimension is four, and all lift-group
elements act cellularly without inversions.

Let \(I\) be the set of vertices fixed by \(C\). Every other cell
has trivial \(C\)-stabilizer, because the fixed locus of its
nontrivial image in \(S_3\) is isolated or empty.
Use \emph{ordinary} cellular cochains with coefficients
\(\mathbb F_p\), so cochains on an infinite set of cells are products,
not compactly supported sums. Apply a free resolution of the finite
cyclic group, of finite rank in each degree, to this cochain complex.
Contractibility identifies the cohomology of the resulting total
complex with \(H^*(C,\mathbb F_p)\).

In the other spectral sequence, every regular cell orbit contributes
a coinduced module, with zero positive-degree group cohomology.
Products have the same property: the resolution has finite rank in
each degree and products are exact. Consequently, for positive
vertical degree \(q\), its only nonzero column is
\[
       E_1^{0,q}=\prod_{i\in I}\mathbb F_p.
 \tag{A.124.9}
\]
For total degree \(n>4\) there is no possible differential from this
column to the row of degree zero or to any other nonzero column.
It follows that
\[
       \prod_{i\in I}\mathbb F_p
          \cong H^n(C,\mathbb F_p)=\mathbb F_p
                    \qquad(n>4).
 \tag{A.124.10}
\]
Thus \(I\) has exactly one element. There is no order-two subgroup,
since a transposition has no fixed point downstairs. Every
order-three subgroup has exactly one fixed point in the cover.

Conversely, every point in (A.124.8) and every chosen lift determine
the unique lift of its three-cycle fixing that point; its cube is
a deck transformation fixing a point and is therefore the identity.
The stabilizer of the chosen lift is exactly \(C_3\). A normalizer
of this subgroup must preserve its unique fixed point, so the
normalizer and centralizer both equal this \(C_3\).
Deck transformations are transitive on each fiber. Hence conjugacy
classes of these subgroups correspond to \(S_3\)-orbits in (A.124.8).
The three-cycles fix each point and a transposition sends \(a\) to
\(-a\), giving four orbits. Since the normalizer cannot invert a
generator, these give eight classes of torsion elements.

\emph{A.124.3. The exact elementary label and completion.}
Abelianizing the marked presentation, the braid relation identifies
\(\sigma\) and \(\tau\); write their common class as \(s\).
The mixed commutator relation gives \(2s=0\).
The two total strand relations then give \(3X=3Y=0\), since their
additional terms are multiples of \(6s\). All other relations become
identities. This proves (A.124.2), including independence of its factors.

Taking \(w=1\) in A.121 gives \(\Pi_0^3=1\), and (A.124.2) gives
\([\Pi_0]=[Y]\ne0\). Orientation-preserving linear homeomorphisms
of \(T\) induce outer automorphisms of the actual configuration
group. They act on the \(C_3^2\) factor by their usual matrices:
the two coordinates record total strand displacement modulo three,
while a diagonal translation changes that displacement by a multiple
of three. Their elementary matrices act transitively on the eight
nonzero vectors of \(\mathbb F_3^2\). We therefore obtain torsion
elements with every nonzero label. There are already exactly eight
classes, so the labels distinguish all of them.

Apply the verified prime-order comparison to obtain the same
classification and centralizers in \(\widehat{\mathcal G}\).
There are no additional finite subgroups: goodness and finite
cohomological dimension make \(\widehat{\mathcal P}\) torsion-free,
so a finite subgroup of \(\widehat{\mathcal G}\) injects into
\(S_3\); the prime-two comparison excludes every subgroup containing
an involution. The remaining nontrivial possibility is \(C_3\).

\emph{A.124.4. The canonical pure transporter.}
Every evaluation of the commutator word \(w\) has trivial image in
the abelianization of its target. Thus \(\Pi_w\) has the same
abelianization \([Y]\) and the same permutation \(\bar r\) as
\(\Pi_0\). A.121 gives its cube one; its nonzero abelian label
makes its order exactly three. The classification supplies a
conjugator \(g\in\widehat{\mathcal G}\) from \(\Pi_0\) to
\(\Pi_w\). Its permutation centralizes \(\bar r\), so lies in
\(\langle\bar r\rangle\). Multiplication on the right by the
appropriate power of \(\Pi_0\) makes \(g\) pure without changing
the conjugacy equation. Two pure choices differ by an element of
\[
 C_{\widehat{\mathcal G}}(\Pi_0)
             \cap\widehat{\mathcal P}=1.
 \tag{A.124.11}
\]
This proves existence and uniqueness in (A.124.3).
The continuous injection
\(p\mapsto p\Pi_0p^{-1}\) from the compact group
\(\widehat{\mathcal P}\) is a homeomorphism onto its closed orbit.
Word evaluation is continuous, and the commutator and pentagon
conditions define a closed set. Composing evaluation with this
inverse proves the asserted continuous selection.

\emph{A.124.5. The exact involution class.}
The actual elliptic section of the original pair
\(\alpha=(-1,\phi)\) preserves the permutation map.
Its formulas in A.122, and the commutator nature of \(\phi\),
give on (A.124.2)
\[
        \Theta(s)=s,\qquad \Theta(X)=-X,
                  \qquad \Theta(Y)=Y.
 \tag{A.124.12}
\]
Thus \(\Theta(\Pi_w)\) has the same torsion label and permutation
as \(\Pi_w\). The preceding argument gives a unique pure
\(h_w\) carrying \(\Pi_w\) to \(\Theta(\Pi_w)\).
Applying \(\Theta\) twice shows that
\(\Theta(h_w)h_w\) centralizes \(\Pi_w\). It is pure, so is one
by (A.124.11) and conjugation. Therefore \(z_w=h_w^{-1}\) satisfies
\(z_w\Theta(z_w)=1\), and
\[
 (z_w\vartheta)\Pi_w(z_w\vartheta)^{-1}=\Pi_w.
 \tag{A.124.13}
\]
The two commuting elements have orders three and two respectively;
the second has nontrivial image in the quotient \(C_2\).
They therefore generate the actual cyclic group of order six in
(A.124.5).

If \(\Pi'=a\Pi_wa^{-1}\) for pure \(a\), uniqueness gives
\[
 h'=\Theta(a)h_wa^{-1},\qquad
 z'=a z_w\Theta(a)^{-1}.
 \tag{A.124.14}
\]
These are exactly the usual nonabelian coboundary changes.
All the words \(\Pi_w\) have already been shown to belong to the
same pure orbit, so their \(C_2\)-classes coincide for fixed
\(\Theta\). Moreover a pure conjugate \(a\Pi_wa^{-1}\) is fixed
by \(\Theta\) precisely when its uniquely normalized transporter
is one, or equivalently when
\(z_w=a^{-1}\Theta(a)\). This proves the final equivalence.

The constructed involution \(z_u\vartheta\) is not asserted to be
pure-conjugate to the original \(\vartheta\). That is the explicit
remaining cohomology condition. Neither the torsion classification
nor this cyclic subgroup supplies the missing boundary equality of
A.122--A.123, identifies the original coherent CM frame, or proves
the two main decisions.
\end{proof}
